\begin{textsample}{POS Dim 3 – 2009\_c6 – Score -2.00 – 2009\_c6/enade\_2009\_admin\_b2a0b87461.txt}  \label{ex:f3_pos_236}
Capítulo 3 Análise Técnica da Prova Este capítulo tem por objetivo apresentar o desempenho dos estudantes de ADMINISTRAÇÃO no ENADE/2009.
Para isso, foram calculadas as estatísticas básicas da prova como um todo, bem como as estatísticas das questões relacionadas à Formação Geral e ao Componente Específico.
Nas Tabelas, são apresentadas as seguintes estatísticas : total da população, de presentes, média, erro-padrão da média, desvio-padrão, mediana, nota máxima, nota mínima e coeficiente de assimetria ( ver definição em 1.4.3 ).
Nas estatísticas das questões discursivas, foram detalhadas também as notas nulas : em branco, por protesto ou por insuficiência.
As estatísticas apresentadas neste capítulo contemplaram, separadamente, os ingressantes, os concluintes e o total de estudantes, e foram calculadas tendo -se em vista as seguintes agregações : ( a ) Região Geográfica e País e ( b ) a Categoria Administrativa e a Organização Acadêmica.
As notas brutas das questões discursivas em cada uma das partes, Formação Geral/Componente Específico, foram calculadas a partir da média simples das notas de cada uma das questões correspondentes.
A primeira questão discursiva da Formação Geral foi desconsiderada e, portanto, a nota bruta referente às questões discursivas de Formação Geral foi igual à nota da questão mantida, a segunda.
No Componente Específico, a nota bruta correspondeu à média das notas nas três questões.
As notas brutas das questões objetivas de cada uma das partes, Formação Geral/Componente Específico, foram calculadas a partir dos acertos dos estudantes considerando -se uma regra de três para a conversão para a escala de 0 a 100, descontando -se aquelas questões que foram desconsideradas pela comissão ou que não passaram pelo teste do índice de discriminação ponto-bisserial ( ver seção 1.4.6 para a descrição do índice ).
As questões com índices de discriminação inferior a 0,2 foram excluídas do cômputo da nota do estudante.
As notas brutas de Formação Geral e do Componente Específico para cada um dos estudantes foram calculadas como médias ponderadas das suas notas nas questões objetivas e discursivas, ponderação diferente para cada parte : C\_FG = 0,6 × C\_FG^OBJ + 0,4 × C\_FG^DIS C\_CE = 0,85 × C\_CE^OBJ + 0,15 × C\_CE^DIS, onde C\_FG é a nota bruta da Formação Geral, C\_FG^OBJ é a nota bruta nas questões objetivas da Formação Geral, C\_FG^DIS é a nota bruta nas questões discursivas da Formação Geral, C\_CE é a nota bruta do Componente Específico, C\_CE^OBJ é a nota bruta nas questões objetivas do Componente Específico, e C\_CE^DIS é a nota bruta nas questões discursivas do Componente Específico.
A nota bruta da prova, C\_GER, por sua vez, é também uma média ponderada da nota de cada uma das partes, Formação Geral/Componente Específico : C\_GER = C\_FG × 0,25 + C\_CE × 0,75. 3.1 Estatísticas Básicas da Prova 3.1.1 Estatísticas Básicas Gerais A Tabela 3.1 apresenta as estatísticas básicas da prova por grupo de estudantes.
O total de estudantes inscritos para a prova foi de 325.351.
Destes, 74,7 \% compareceram, sendo que a presença foi maior entre concluintes ( 87,3 \% ) do que entre ingressantes ( 67,1 \% ).
A média geral da prova foi de 35,06, sendo que os ingressantes obtiveram média mais baixa, de 31,88, que os concluintes, cuja média foi 39,12.
O desvio-padrão geral foi de 14,3, sendo o do grupo de ingressantes menor ( 12,9 ), do que o do grupo dos concluintes ( 15,0 ), aqueles possuíram uma distribuição mais homogênea nas notas.
A nota máxima foi 97,7, obtida por pelo menos um concluinte ao passo que a maior nota obtida por um ingressante foi 92,8.
Os coeficientes de assimetria para os dois grupos considerados, o de concluintes e o de ingressantes, bem como para a população como um todo foram todos positivos, indicando uma maior concentração à direita da distribuição e um maior espalhamento à esquerda.
Uma comparação entre o desempenho de ingressantes e concluintes pode ser realizada a partir dos dados apresentados no Gráfico 3.1.
Cumpre destacar que as notas dos concluintes, de um modo geral, foram mais elevadas do que as dos ingressantes, cuja distribuição de notas está à esquerda dos primeiros.
A ideia de melhor desempenho foi reforçada pelas notas medianas dos dois grupos : 37,8 entre os concluintes e 30,4 entre os ingressantes.
Nas faixas de notas até 40 pontos o percentual de ingressantes foi maior que o dos concluintes, situação esta que se inverte para as faixas de notas superiores a 40 pontos.
Constatou -se que aproximadamente 9,4 \% dos ingressantes e 23,5 \% dos concluintes situaram -se nas faixas de notas superiores a 50,0 pontos.
O Gráfico 3.2 apresenta informações referentes ao desempenho geral de ingressantes e concluintes, comparando os resultados considerando a Região, Categoria Administrativa e Organização Acadêmica.
Levando -se em conta as notas médias dos estudantes em cada Região, observa -se que, nos dois grupos, há uma variação expressiva entre as Regiões.
As notas médias mais elevadas entre os concluintes ocorreram nas Regiões Sul ( 42,09 ) e Sudeste ( 39,19 ), ambas acima da média nacional ( 39,12 ).
Em relação aos ingressantes, as médias mais altas foram obtidas nas Regiões Sul ( 32,80 ) e Sudeste ( 32,01 ) nas demais Regiões a média é inferior à média nacional ( 31,87 ).
Na Região Centro-Oeste foi obtida a média de notas mais baixa entre os ingressantes ( 29,91 ).
Entre os concluintes, a nota mais baixa foi verificada na Região Norte ( 35,35 ).
Cursos nesta área foram oferecidos em instituições federais, estaduais, municipais e privadas.
Levando -se em conta os agrupamentos dos estudantes em categorias administrativas, observou -se que a pontuação média mais elevada entre os concluintes e ingressantes foi encontrada em instituições de origem federal ( respectivamente 53,80 e 42,25 ) – ambas acima da média nacional.
Tendo como foco as notas médias dos estudantes agrupados por Organização Acadêmica, notamos que para os ingressantes, os Centros Federais de Educação Tecnológica apresentaram as maiores médias ( 47,17 ), ao passo que as piores notas médias foram observadas nas Faculdades ( 30,66 ).
Entre os concluintes, as maiores médias foram alcançadas pelos estudantes dos Institutos Federais de Educação, Ciência e Tecnologia ( 51,18 ), enquanto as piores foram, também obtidas pelos estudantes de Faculdades ( 37,53 ). 3.1.2 Estatísticas Básicas em Formação Geral A Tabela 3.2 apresenta as estatísticas básicas de Formação Geral por grupo de estudantes.
O total de estudantes inscritos para a prova foi de 325.351.
Destes, 74,7 \% compareceram, sendo que a presença foi maior entre concluintes ( 87,3 \% ) do que entre ingressantes ( 67,1 \% ).
A média geral da parte da prova de Formação Geral foi de 41,50, sendo que os ingressantes obtiveram média mais baixa, 39,38, que os concluintes, cuja média foi 44,22.
O desvio-padrão geral foi de 20,5, sendo o do grupo de ingressantes ( 20,1 ), e o dos concluintes ( 20,7 ), aqueles possuíram uma distribuição mais homogênea nas notas.
A nota máxima foi 100, obtida por estudantes de ambos os grupos.
Para esta parte, assim como para a nota geral, os coeficientes de assimetria foram positivos para todos os grupos.
Como, porém, os valores estavam bem perto de zero, estes indicaram uma ligeira inclinação para a direita.
O Gráfico 3.3 permite a comparação entre concluintes e ingressantes relativamente ao desempenho na parte da prova que aborda a Formação Geral.
Nesta parte, se excluídas as provas com nota zero, tal como na prova como um todo, também os concluintes apresentaram desempenho superior, embora seja perceptível um maior equilíbrio.
A mediana dos concluintes é superior à dos ingressantes, assim como a média.
Dos ingressantes, aproximadamente 29,7 \% alcançaram nota superior a 50 ; dentre os concluintes, o percentual dos estudantes nesta faixa de notas foi muito semelhante e igual a 38,7 \%.
No Gráfico 3.4, são apresentadas as informações referentes ao desempenho dos ingressantes e dos concluintes na parte de Formação Geral, em diferentes agregações como Região, Categoria Administrativa e Organização Acadêmica.
Considerando -se as notas médias dos estudantes segundo cada Região, observou -se que, tanto no grupo dos ingressantes quanto no grupo dos concluintes, houve diferenças marcantes.
As médias variaram, entre os concluintes, de 40,03 na Região Norte a 48,11 na Região Sul.
Entre os ingressantes, os valores foram de 35,82 na Região Norte a 41,16 na Região Sul.
Em todas as Regiões, os concluintes apresentaram médias maiores do que os ingressantes.
Tendo como foco as categorias administrativas, destacou -se que todas as esferas estavam representadas, sendo que as notas médias das instituições federais foram superiores para concluintes ( 55,84 ), enquanto as notas médias das instituições estaduais foram superiores para ingressantes ( 50,40 ).
Considerando -se o tipo de Organização Acadêmica, Centros Federais de Educação Tecnológica apresentaram os maiores valores, 53,81 para concluintes e 58,91 para ingressantes.
Como podemos observar, tanto os Centros Federais de Educação Tecnológica quanto os Institutos Federais de Educação apresentaram notas médias maiores para ingressantes do que para concluintes. 3.1.3 Estatísticas Básicas do Componente Específico A Tabela 3.3 apresenta as estatísticas básicas referentes ao Componente Específico da prova, tendo sua média geral, de 32,97 inferior ao desempenho da prova de ADMINISTRAÇÃO como um todo, cuja média foi 35,06, como apresentado na Tabela 3.1, que por sua vez é inferior ao desempenho na Formação Geral, que teve média igual a 41,50, conforme mostrado na Tabela 3.2.
Os concluintes apresentaram um desempenho médio ( 37,48 ) superior ao dos ingressantes ( 29,43 ), o mesmo ocorrido para a nota da Formação Geral, havendo ainda indícios de maior homogeneidade entre as notas médias dos ingressantes do que entre as dos concluintes, como sugeriu o desvio padrão, menor no grupo dos ingressantes do que no dos concluintes.
A nota máxima obtida por concluintes ( 97,0 ) foi maior que a dos ingressantes ( 92,2 ).
As distribuições de notas para cada grupo se apresentaram com uma inclinação para notas maiores, como evidenciados pelos coeficientes de assimetria positivos.
A situação de assimetria foi confirmada pelo Gráfico 3.5 que proporciona uma comparação do desempenho entre concluintes e ingressantes em relação ao Componente Específico da prova.
Claramente, a distribuição das notas médias dos ingressantes, à direita daquela dos concluintes, evidenciou o percentual mais alto de notas médias nas faixas mais baixas para este grupo.
As medianas das distribuições das notas dos ingressantes e dos concluintes foram 27,4 e 35,7, respectivamente, enfatizando o melhor desempenho dos concluintes nesta parte da prova.
Assim como os Gráficos 3.2 e 3.4, o Gráfico 3.6 apresenta uma comparação dos resultados em relação às Regiões, Categoria Administrativa e Organização Acadêmica, agora levando em conta o desempenho dos ingressantes e dos concluintes no Componente Específico da prova.
Tendo como foco as notas médias dos estudantes, em determinada Região, observa -se que, entre os concluintes, as mais elevadas foram encontradas nas Regiões Sul ( 40,13 ) e Sudeste ( 37,33 ), valores acima da média nacional para esta parte ( 37,47 ).
Entre os ingressantes, as maiores notas ocorreram nas Regiões Sul ( 30,07 ) e Nordeste ( 29,96 ) todos acima da média nacional ( 29,43 ).
Na análise, segundo a Categoria Administrativa verificou -se a presença de todas as modalidades : entidades federais, estaduais, municipais e privadas.
As médias mais elevadas, tanto entre os concluintes como entre os ingressantes, foram encontradas em instituições de origem federal ( 53,16 entre os concluintes e 39,78 entre os ingressantes ), enquanto nas instituições estaduais as médias foram, respectivamente, 43,20 e 37,54, para concluintes e ingressantes.
Já nas instituições municipais, as médias foram as menores entre as categorias administrativas, tanto para concluintes como para ingressantes, respectivamente 35,81 e 28,10.
Quanto à Organização Acadêmica, as universidades, centros universitários, faculdades, Centro Federal de Educação Tecnológica e Instituto Federal de Educação, Ciência e Tecnologia foram representadas neste curso.
Considerando -se o tipo de Organização Acadêmica, os Centros Federais de Educação Tecnológica apresentaram os maiores valores, 46,07 para concluintes e 43,31 para ingressantes. 3.2 Análise das Questões Objetivas 3.2.1 Formação Geral A Tabela 3.4 apresenta as estatísticas básicas relativas às oito questões objetivas da Formação Geral dos estudantes.
Adiantamos que a primeira questão foi desconsiderada pelo critério da correlação ponto-bisserial.
Como pode ser observado na Tabela 3.4, os estudantes concluintes obtiveram desempenho superior, com nota média de 53,94, em relação aos estudantes ingressantes, cuja nota média foi 49,69.
A variabilidade entre os concluintes foi cerca de 1 \% menor do que entre os ingressantes, como indicaram os desvios-padrão de 21,9, para o grupo de ingressantes e 22,2 para o grupo dos concluintes.
A amplitude total das notas foi a mesma para ambos os grupos : nota mínima de 0 e máxima de 100.
A mediana foi 57,1 para os concluintes e 42,8 para os ingressantes.
A assimetria, medida pelo coeficiente de assimetria ( skewness ) foi um pouco maior em módulo para os ingressantes ( 0 ) do que para os concluintes ( -0,1 ), ainda que ambos fossem muito similares, este coeficiente indicou que a distribuição de notas dos ingressantes foi mais simétrica do que a dos concluintes.
A Tabela 3.5 apresenta a classificação das questões objetivas da Formação Geral, segundo o índice de facilidade ( questões em vermelho foram desconsideradas pelo critério ponto-bisserial ).
Como as questões foram comuns a todas as áreas, a tabela apresenta a classificação considerando -se os concluintes de todas as áreas e aqueles do curso de ADMINISTRAÇÃO.
De acordo com os índices obtidos, as questões objetivas da prova foram assim avaliadas para a área de ADMINISTRAÇÃO : das oito questões, duas tiveram entre 60 \%, inclusive, e 85 \% de acertos e foram classificadas como fáceis ( somente uma A Tabela 3.6 descreve os resultados das análises das questões objetivas relativas à Formação Geral, segundo o poder de discriminação, utilizando -se para tal o índice de discriminação ponto-bisserial.
Como as questões foram comuns a todas as áreas, a tabela apresenta a classificação considerando -se os concluintes de todas as áreas e aqueles do curso de ADMINISTRAÇÃO.
A desconsideração da questão, segundo este critério, considerou o universo de todos os concluintes.
Nenhuma das oito questões apresentou índices igual ou acima de 0,40 ( nenhuma para a população de concluintes como um todo ), sendo, portanto, classificadas como muito boas nesta característica, para este grupo de estudantes, duas foram classificadas como boa com índice entre 0,30, inclusive, e 0,40 ( três para os concluintes como um todo ) e cinco foram classificadas como médias, com índice entre 0,20, inclusive, e 0,30 ( quatro para os concluintes como um todo ).
Uma das questões foi classificada como fraca tanto para a população de concluintes como um todo como para a de concluintes de ADMINISTRAÇÃO e foi desconsiderada pelo critério. para o total de concluintes ) ; outras três foram consideradas medianas, alcançando entre 40 \%, inclusive, e 60 \% de acertos ( três para o total de concluintes ) ; duas obtiveram entre 15 \%, inclusive, e 40 \% de acertos, sendo consideradas difíceis ( três para o total de concluintes ).
Por fim, uma questão foi classificada como muito difícil tendo obtido percentual menor do que 15 \% de acertos ( uma para o total de concluintes ), coincidentemente a questão desconsiderada pelo critério ponto-bisserial.
Nenhuma questão foi classificada na categoria muito fácil.
Em linhas gerais podemos dizer que os concluintes desta área tiveram um melhor desempenho nesta parte do que o total da população de concluintes de todos os cursos avaliados em 2009.
A Tabela 3.6 descreve os resultados das análises das questões objetivas relativas à Formação Geral, segundo o poder de discriminação, utilizando -se para tal o índice de discriminação ponto-bisserial.
Como as questões foram comuns a todas as áreas, a tabela apresenta a classificação considerando -se os concluintes de todas as áreas e aqueles do curso de ADMINISTRAÇÃO.
A desconsideração da questão, segundo este critério, considerou o universo de todos os concluintes.
Nenhuma das oito questões apresentou índices igual ou acima de 0,40 ( nenhuma para a população de concluintes como um todo ), sendo, portanto, classificadas como muito boas nesta característica, para este grupo de estudantes, duas foram classificadas como boa com índice entre 0,30, inclusive, e 0,40 ( três para os concluintes como um todo ) e cinco foram classificadas como médias, com índice entre 0,20, inclusive, e 0,30 ( quatro para os concluintes como um todo ).
Uma das questões foi classificada como fraca tanto para a população de concluintes como um todo como para a de concluintes de ADMINISTRAÇÃO e foi desconsiderada pelo critério.
Na sequência, no Quadro 3.1, estão especificadas, por questão, as habilidades aferidas pelas questões de Formação Geral da prova da ADMINISTRAÇÃO, bem como os resultados da análise relativa aos índices de facilidade e de discriminação de cada questão.
O índice de facilidade para a área de ADMINISTRAÇÃO teve uma variação considerável, oscilando entre 0,10 até 0,81 ( na população de concluintes o intervalo foi de 0,08 a 0,75 ) e o de discriminação ficou entre 0,02 e 0,39 ( na população de concluintes o intervalo foi de 0,02 a 0,39 ).
As questões 6 e 7, que tiveram em comum aferir as habilidades de “ Analisar e criticar informações ” ; ” Extrair conclusões por indução e/ou dedução ”, ” Estabelecer relações, comparações e contrastes em diferentes situações ”, “ Detectar contradições ”, “ Fazer escolhas valorativas avaliando consequências ”, “ Questionar a realidade ” e “ Argumentar coerentemente ” foram as que apresentaram, para a população de concluintes de ADMINISTRAÇÃO, maior poder discriminatório com índice de 0,37 e 0,39, respectivamente.
A questão mais fácil, com 81 \% de acertos, foi a de número 8, que aferiu as habilidades de “ Ler e interpretar textos ”, “ Analisar e criticar informações ”, “ Extrair conclusões por indução e/ou dedução ”, “ Estabelecer relações, comparações e contrastes em diferentes situações ”, “ Detectar contradições ”, “ Fazer escolhas valorativas avaliando consequências ”, “ Questionar a realidade ” e “ Argumentar coerentemente ”.
De outra parte, a questão com menor índice de facilidade e de discriminação em Formação Geral foi a questão 1, que foi desconsiderada, e que deveria aferir também as mesmas habilidades listadas anteriormente.
O Gráfico 3.7 mostra, a título de ilustração, o comportamento das respostas à questão 2 de Formação Geral para os estudantes concluintes da área de ADMINISTRAÇÃO como função da nota geral na prova.
Trata -se de uma questão considerada, com respeito à dificuldade, de classificação médio para os concluintes de ADMINISTRAÇÃO e também para o conjunto de concluintes de todas as áreas.
O poder de discriminação foi também classificado como “ médio ” tanto para os estudantes da área quanto para o conjunto de todos os concluintes.
Neste Gráfico, cada uma das cinco curvas, representou o percentual de respostas em determinada alternativa da questão, em função do escore total dos estudantes da área de ADMINISTRAÇÃO.
A curva vermelha correspondeu à alternativa C, a correta para esta questão.
Assim, observou -se que entre os estudantes com notas mais baixas – desempenho geral mais fraco na prova como um todo –, a tendência foi a escolha de alternativas incorretas, especialmente a alternativa D – curva verde escura, seguida da alternativa E – curva em roxo.
A proporção de estudantes que selecionaram a resposta correta C, aumentou com o desempenho deste grupo de estudantes chegando ao acerto de 90 \% dos estudantes com notas pelo menos 75.
Esta análise permitiu verificar como a questão discriminou entre os grupos de desempenho.
Os Gráficos relativos às demais questões de Formação Geral constam do Anexo I. 3.2.2 Componente Específico A Tabela 3.7 apresenta as estatísticas básicas em relação às questões objetivas do Componente Específico.
Como pode ser observado, os estudantes concluintes obtiveram desempenho superior, com uma média de 40,67 pontos, em relação aos estudantes ingressantes com uma média de 32,24 pontos.
Os desvios-padrão indicaram que a variabilidade entre os concluintes foi maior ( 17,5 ) do que entre os ingressantes ( 14,9 ).
Os coeficientes de assimetria para os dois grupos foram positivos indicando uma inclinação para a direita.
A assimetria foi, porém, maior entre os concluintes.
A análise das questões objetivas do Componente Específico em ADMINISTRAÇÃO, em relação ao grau de facilidade da prova ( ver Tabela 3.8 ), mostrou que, das 27 questões, 17 encontraram -se na categoria difícil ( entre 15 \% e 40 \%, não incluso, de acertos ), 6 foram de nível médio ( entre 40 \% e 60 \%, não incluso, de acertos ), 2 estavam na categoria fácil ( entre 60 \% e 85 \%, não incluso, de acertos ) e nenhuma obteve a classificação muito fácil ( a partir de 85 \% de acertos ).
Nenhuma questão obteve percentual de acertos menor do que 15 \% ( classificação muito difícil ), 2 questões foram desconsideradas pela Comissão Assessora da Área de ADMINISTRAÇÃO.
Das 27 questões objetivas relativas ao Componente Específico ( ver Tabela 3.9 ), 5 obtiveram o índice de discriminação muito bom, 5 obtiveram o índice de discriminação bom e 8 tiveram nível de discriminação médio.
O índice fraco que eliminaria a questão, foi alcançado por 7 das questões, sendo que 2 questões foram desconsideradas pela comissão.
Na sequência, no Quadro 3.2 estão especificadas, por questão, as habilidades aferidas pelas questões do Componente Específico da prova, bem como os resultados da análise relativa aos índices de facilidade e de discriminação de cada questão.
O Gráfico 3.8 mostra, a título de ilustração, o comportamento da questão 13 do Componente Específico para os estudantes concluintes da área de ADMINISTRAÇÃO como função da nota geral na prova.
Trata -se de uma questão considerada média para os concluintes de ADMINISTRAÇÃO.
O poder de discriminação foi bom para os estudantes da área.
Neste Gráfico, cada uma das cinco curvas, representou o percentual de respostas em determinada alternativa da questão, em função da nota total dos estudantes concluintes de ADMINISTRAÇÃO.
A curva vermelha correspondeu à alternativa C, a correta para esta questão.
Assim, observou -se que entre os estudantes com escores mais baixos – desempenho geral mais fraco na prova como um todo –, a tendência foi a escolha de uma alternativa incorreta, a alternativa E – curva em roxo.
A proporção de estudantes que selecionaram a resposta correta aumentou com o desempenho deste grupo de estudantes chegando ao acerto de quase 100 \% dos estudantes com notas pelo menos 87,5.
Esta análise permitiu verificar como a questão discriminou entre os grupos de desempenho.
Os Gráficos relativos às demais questões do Componente Específico constam do Anexo I. 3.3 Análise das Questões Discursivas 3.3.1 Formação Geral A Tabela 3.10 apresenta as estatísticas básicas da prova por grupo de estudantes.
A análise dos resultados de desempenho dos estudantes nas questões discursivas que são apresentados na tabela 3.10 e no gráfico 3.9, mostrou que as notas foram mais baixas no conjunto dessas questões que no das questões objetivas.
O total de estudantes inscritos para a prova foi de 325.351.
Destes, 74,7 \% compareceram, sendo que a presença foi maior entre concluintes ( 87,3 \% ) do que entre ingressantes ( 67,1 \% ).
A média geral da prova foi de 26,53, sendo que os ingressantes obtiveram média mais baixa, de 24,01, que os concluintes, cuja média foi 29,74.
O desvio-padrão geral foi de 33,3, sendo o do grupo de ingressantes menor ( 32,2 ), em relação ao grupo dos concluintes ( 34,5 ), aqueles possuíram uma distribuição mais homogênea nas notas.
A nota máxima foi 100,0, atingida por concluintes e ingressantes.
Os coeficientes de assimetria para os dois grupos considerados, o de concluintes e o de ingressantes, bem como para a população como um todo foram todos positivos, indicando uma maior concentração à esquerda da distribuição.
Para todos os grupos, cerca de metade dos estudantes presentes obtiveram nota zero nas questões discursivas da Formação Geral.
Das 134.911 provas com nota zero, 5.828 ( 4,3 \% ) continham um protesto, 34.201 ( 25,4 \% ) estavam em branco e 94.882 ( 70,3 \% ) foram respondidas, porém de forma inadequada.
Entre os concluintes, a minoria ( 22,7 \% ) dentre aqueles com nota zero estavam com a prova em branco.
Entre os ingressantes, 68,8 \% das notas nulas ( zero ) decorreram do desempenho do próprio estudante verificadas na correção.
Uma comparação entre o desempenho de ingressantes e concluintes pode ser realizada a partir dos dados apresentados no Gráfico 3.9.
As notas dos concluintes, de um modo geral, foram ligeiramente mais elevadas do que as dos ingressantes, cuja distribuição de notas está à direita dos primeiros.
A ideia de melhor desempenho foi reforçada pela proporção de notas abaixo de 10 : 51,6 entre os concluintes e 59,0 entre os ingressantes.
Na sequência, discorrer -se-á, especificamente, sobre uma das questões discursivas de Formação Geral, a número 10 já que a número 9 foi desconsiderada.
Em especial, serão analisados os conteúdos englobados na questão e os resultados obtidos pelos estudantes.
Ao final, serão apresentados os comentários da banca corretora a respeito do observado nas respostas dos estudantes, suas impressões e conclusões.
Cumpre esclarecer que, tendo em vista que as questões discursivas de Formação Geral foram padronizadas, ou seja, constam de todas as provas, os comentários da banca foram os mesmos para todas as áreas acadêmicas, eis que direcionados a todos os estudantes que participaram do ENADE/2009. 3.3.1.1 Análise da Questão Discursiva 9 de Formação Geral DESCONSIDERADA 3.3.1.2 Análise da Questão Discursiva 10 de Formação Geral A Questão Discursiva nº. 10 buscou avaliar as habilidades : • Ler e interpretar textos ; • Analisar e criticar informações ; • Extrair conclusões por indução e/ou dedução ; • Estabelecer relações, comparações e contrastes em diferentes situações ; • Fazer escolhas valorativas avaliando consequências ; • Questionar a realidade ; • Argumentar coerentemente.
O Quadro 3.3 apresenta as habilidades/competências, temas predominantes e estatísticas básicas da questão 10 de Formação Geral por grupo de estudantes.
Cumpriu notar que como a questão 9 foi desconsiderada, a nota referente às questões discursivas da Formação Geral é idêntica à nota da questão 10.
A média geral desta questão foi de 26,53, sendo que os ingressantes obtiveram média mais baixa, de 24,01, que os concluintes, cuja média foi 29,74.
O desvio-padrão geral foi de 33,3, sendo que o do grupo de ingressantes foi ligeiramente menor ( 32,2 ), em relação ao grupo dos concluintes ( 34,5 ), aqueles possuíram uma distribuição mais homogênea nas notas.
A nota máxima foi 100,0, atingida por concluintes e ingressantes.
Os coeficientes de assimetria para os dois grupos considerados, o de concluintes e o de ingressantes, bem como para a população como um todo foram positivos, indicando uma maior concentração à esquerda da distribuição.
Para todos os grupos, cerca de metade dos estudantes presentes obtiveram nota zero nesta questão discursiva de Formação Geral.
Das 134.911 provas com nota zero, 5.828 ( 4,3 \% ) continham um protesto, 34.201 ( 25,4 \% ) estavam em branco e 94.882 ( 70,3 \% ) foram respondidas, porém de forma inadequada.
Entre os concluintes, a minoria ( 22,7 \% ) dentre aqueles com nota zero estava com a prova em branco.
Entre os ingressantes, 68,8 \% das notas nulas ( zero ) decorreram do desempenho do próprio estudante verificadas na correção.
Uma comparação entre o desempenho de ingressantes e concluintes para esta questão foi exatamente igual a já realizada para a nota discursiva em Formação Geral a partir dos dados apresentados no Gráfico 3.9. 3.3.1.3 Comentários sobre as respostas à Questão Discursiva 10 Apresentação A questão 10 de Formação Geral apresentou um questionamento capaz de analisar a habilidade do estudante em reconhecer o que é um direito social, optar por um direito social vigente que possa diminuir as desigualdades de renda familiar no Brasil e discutir os argumentos em favor do direito social escolhido como suporte à resposta.
Para análise das respostas apresentadas pelos estudantes, houve uma leitura inicial bastante criteriosa.
Foram considerados para efeito de pontuação, os seguintes argumentos : • Acesso à educação pública, gratuita e de qualidade, garantindo ao cidadão melhores condições de vida ; • Assistência à saúde, em seu contexto mais amplo, garantindo os cuidados essenciais de saúde, baseados em técnicas apropriadas ao alcance de todas as pessoas, contando com a participação da população ; • Permanência do estudante na escola, em todos os níveis escolares ( da educação infantil à educação superior, possibilitando a qualificação profissional ), o que garante melhores condições de trabalho e remuneração.
Considerou -se a qualificação e o desempenho docente, desenvolvendo nos estudantes melhores expectativas de aprendizagem e interesse pelo conteúdo da escola.
Os professores desenvolvendo metodologias adequadas e conhecimento pleno quanto a questões de aprendizagem ; • Condições dignas de trabalho, com remuneração que garanta qualidade de vida.
Foram consideradas as condições de trabalho, a preparação profissional para o trabalho e os cursos de aperfeiçoamento, oportunizando momentos de aprendizagem e desenvolvimento, a partir de situações desencadeadoras ; • Ser proprietário do imóvel em que reside, reduzindo os gastos com aluguel, garantindo melhores condições de vida, equilíbrio financeiro e construção do bem-estar familiar.
Desempenho dos estudantes na questão É importante salientar que havia respostas que desenvolveram adequadamente o questionamento apresentado, com fundamentação baseada em argumentos coerentes.
Contudo, em um número considerável de respostas, havia apenas um argumento e não dois, conforme fora solicitado.
Muitos estudantes informaram sobre problemas da educação, de um modo geral.
Algumas observações mencionadas pelos corretores durante o processo de correção foram : • há preocupação com uma educação de qualidade, com sugestões de cursos profissionalizantes e métodos preparatórios de aulas. • a educação é revelada como de baixa qualidade, exaltando -se treinamentos e cursos de atualização para a classe docente. • a saúde como essência para a vida.
O SUS é colocado como parte integrante de direito social, garantindo às pessoas fácil acesso. • oportunidade de emprego como forma de ascensão social, direito, cidadania. • os políticos brasileiros como causadores do despreparo educacional.
As respostas analisadas apresentaram norma culta, coesão e coerência razoáveis, que garantiram a compreensão da mensagem desenvolvida pelos estudantes.
Os erros ortográficos foram poucos e não comprometeram a discussão sobre o assunto e sobre os pontos de vista abordados.
Quanto à coerência, houve um desvio em relação ao objetivo da questão, já que alguns estudantes, através dos seus argumentos retrataram reivindicações, reclamações e situações de problemas enfrentados pelos brasileiros, citando golpes políticos.
Houve erros quanto à concordância verbal e nominal.
Houve, também, dificuldade em utilizar os conectores corretamente.
Alguns estudantes apresentaram em suas respostas redundância e prolixidade.
O desempenho foi considerado satisfatório, pois a maioria dos estudantes soube responder ao questionamento.
Houve domínio do conteúdo, pois as ideias, os argumentos e a originalidade corresponderam ao esperado. 3.3.2 Componente Específico Nas questões discursivas do Componente Específico – questões de números 38, 39 e 40 –, a diferença de desempenho entre os grupos não foi muito acentuada, inclusive menor do que para as questões discursivas de Formação Geral.
A Tabela 3.11 contém as estatísticas da média destas questões.
Observou -se que no grupo dos concluintes, a nota média alcançada foi 19,71, enquanto que no grupo dos ingressantes foi de 13,83, ambas menores do que as obtidas nas questões discursivas de Formação Geral, que foram de 24,01 para os ingressantes e 29,74 para os concluintes ( Tabela 3.10 ).
As médias foram também menores dos que as obtidas para as questões objetivas do Componente Específico ( Tabela 3.7 ) : respectivamente 32,24 e 40,67 para ingressantes e concluintes.
A análise de cada uma destas questões será feita nas subseções seguintes.
Os valores positivos para a medida de assimetria indicaram uma distribuição com concentração à direita e espalhamento à esquerda.
Dos 242.915 presentes, 31,3 \% obtiveram nota zero.
A tabela 3.12 apresenta um detalhamento da situação destes estudantes.
Do total de ingressantes e concluintes, 42,7 \% não fizeram nenhuma das 3 questões.
Esta proporção foi ligeiramente menor entre os concluintes ( 38,2 \% ) do que entre os ingressantes ( 45,4 \% ).
A combinação de questões em branco e preenchida correspondeu a mais de 95 \% das provas com nota nula.
Menos de 1 \% no total destas provas corresponderam a 3 questões com protesto.
O Gráfico 3.10 apresenta a distribuição das notas nas questões discursivas do Componente Específico por grupo de estudantes.
A moda se situou no primeiro intervalo, aquele com valores entre 0 e 10 incluindo os extremos do intervalo.
A impressão de assimetria informada pelos coeficientes de assimetria positivos ( Tabela 3.11 ) foi reforçada no gráfico pela concentração no primeiro intervalo. 3.3.2.1 Análise da Questão Discursiva 38 do Componente Específico A questão 38, cujas habilidades foram aferidas encontra -se descrita no Quadro 3.4, que ainda contém outras informações, seguido da Tabela 3.13 e do Gráfico 3.11, foi aquela em que os grupos de estudantes observados obtiveram um desempenho intermediário entre as questões discursivas do Componente Específico.
A nota média do grupo dos ingressantes foi de 13,73, e a dos concluintes 19,88.
O percentual de notas zero obtidas nos dois grupos foi de 55,6 \%, maior entre os ingressantes ( 60.6 \% ) do que entre os concluintes ( 49,2 \% ).
Um pouco mais da metade dos ingressantes ( 54,3 \% ) e 45,6 \% dos concluintes que tiraram zero deixaram a questão em branco ( ver Tabela 3.13 ).
Uma pequena minoria ( 3,7 \% ) utilizou a questão para protestar.
Apesar do baixo desempenho, em ambos os grupos a nota máxima de 100,0 pontos foi alcançada. 3.3.2.2 Comentários sobre as respostas à Questão Discursiva 38 A questão apresenta como objetivo avaliar a capacidade que o estudante tem de efetuar análise financeira de um investimento e do seu impacto, utilizando os conhecimentos da Contabilidade Financeira.
Trata -se de uma questão que requer maior conhecimento específico por parte dos estudantes.
Os estudantes do primeiro ano do curso de Administração encontraram muitas dificuldades em responder de forma correta à questão.
Apresenta a questão uma boa O Gráfico 3.11 apresenta a distribuição das notas na questão 38 discursiva do Componente Específico por grupo de estudantes.
A moda se situou no primeiro intervalo, aquele com valores entre 0 e 10 incluindo os extremos do intervalo.
A impressão de assimetria informada pelos coeficientes de assimetria positivos ( Quadro 3.4 ) foi reforçada no gráfico pela concentração no primeiro intervalo.
A distribuição apresentou outros máximos locais no quinto e no sétimo intervalos, para concluintes. 3.3.2.2 Comentários sobre as respostas à Questão Discursiva 38 A questão apresenta como objetivo avaliar a capacidade que o estudante tem de efetuar análise financeira de um investimento e do seu impacto, utilizando os conhecimentos da Contabilidade Financeira.
Trata -se de uma questão que requer maior conhecimento específico por parte dos estudantes.
Os estudantes do primeiro ano do curso de Administração encontraram muitas dificuldades em responder de forma correta à questão.
Apresenta a questão uma boa contextualização no seu enunciado e devido ao conteúdo e ao conhecimento exigido para respondê -la, considera -se a questão de nível de dificuldade alto.
Quanto ao desempenho dos estudantes, notou -se que a maioria das respostas demonstrou que os mesmos não possuíam domínio adequado do conteúdo para responder à questão.
Envolvia conhecimento na área de gestão financeira e análise de projetos.
Os estudantes utilizaram de forma desmedida, fora do contexto, os termos técnicos de Contabilidade, como por exemplo : vendas, despesas, depreciação, opinião subjetiva do problema social de desemprego.
A maioria dos estudantes fixou em aumento de produção e não conseguiu ter uma visão ampla dos impactos que a pergunta retratava.
Numa análise mais detalhada e considerando que a questão apresentava dois itens para análise e resposta, verificamos que : Item 1 – as respostas foram evasivas e os estudantes, muitas vezes, reafirmaram que deveria ser feita uma análise “ custo/benefício. ” Item 2 – um número grande de estudantes identificou com mais clareza alguns impactos na produção, como por exemplo, os ganhos de qualidade e eficiência, com a ressalva de que o enunciado da questão facilitava o raciocínio.
Já os impactos nas áreas financeiras e recursos humanos careceram de conhecimento específico por parte da grande maioria dos estudantes que procurava quase sempre utilizar uma linha de raciocínio muito simplória, diferente do que a pergunta solicitava.
Os erros mais comuns encontrados foram respostas evasivas, com repetição do enunciado e fuga ao tema, apesar dos questionamentos terem sido bem claros.
Observou -se também que os estudantes, apesar da aparente falta de domínio do conteúdo, tentaram esboçar alguma coerência e coesão nas respostas, tentando acompanhar informações citadas no enunciado da questão.
Também, se observou o esforço para a utilização da norma culta no processo de comunicação por escrito, nem sempre atingido com sucesso.
Percebeu -se que nessa questão, os estudantes tiveram dificuldades de efetuar análise financeira para cálculo da TIR e mais facilidade para responder sobre os impactos, devido à visualização mental do contexto. 3.3.2.3 Análise da Questão Discursiva 39 do Componente Específico A questão 39, cujas habilidades aferidas encontra -se descrita no Quadro 3.5, que ainda contém outras informações, seguido da Tabela 3.14 e do Gráfico 3.12, foi aquela em que os grupos de estudantes observados obtiveram melhor desempenho entre as questões discursivas do Componente Específico.
A nota média do grupo dos ingressantes foi de 16,54 e a dos concluintes 20,74.
Isto ocorre, principalmente pelo grande percentual de notas nulas : 54,5 \% entre os ingressantes e 46,9 \% entre os concluintes.
Das provas com notas nulas nesta questão, uma significativa proporção foi deixada em branco : 49,2 \% entre os ingressantes e 42,4 \% entre os concluintes ( ver Tabela 3.14 ).
Entre os com nota nula, uma pequena minoria ( 2,4 \% ) utilizou a questão para protestar.
Apesar do baixo desempenho, em ambos os grupos a nota máxima de 100,0 pontos foi alcançada.
Os coeficientes de assimetria positivos indicaram uma inclinação para a esquerda e cauda para a direita.
O Gráfico 3.12 apresenta a distribuição das notas na questão 39 discursiva do Componente Específico por grupo de estudantes.
A moda se situou no primeiro intervalo, aquele com valores entre 0 e 10 incluindo os extremos do intervalo.
A impressão de assimetria informada pelos coeficientes de assimetria positivos ( Quadro 3.5 ) foi reforçada no gráfico pela concentração no primeiro intervalo.
A distribuição apresentou um máximo local no quinto intervalo, de notas entre 40 e 50. 3.3.2.4 Comentários sobre as respostas à Questão Discursiva 39 A questão tem como objetivo avaliar o estudante quanto ao seu conhecimento sobre a aplicabilidade da “ Teoria da Liderança Contingencial ”, como suporte para resolução de um problema de uma empresa contextualizada.
O contexto no enunciado da questão também foi desenvolvido de forma que pudesse \textbf{ajudar} o estudante a desenvolver seus argumentos baseados na liderança contingencial.
O nível de dificuldade da questão é baixo, tratou -se de um conteúdo discutido desde os primeiros períodos do curso de Administração.
Os estudantes, mesmo sem recordar o conceito da Teoria de Liderança Contingencial, conseguiram se transportar para a situação relatada e responderam de forma simples ao que a pergunta questionava.
Notou -se que os erros mais comumente encontrados foram a falta de objetividade e coerência, fugindo do tema e utilizando “ frases prontas ” do enunciado.
Alguns estudantes não conseguiram concluir a análise de todos os blocos, comumente comentava -se o investimento apenas para um bloco e tiveram dificuldades de responder à questão de forma coesa e coerente.
Explicitando, faltou mais coerência nas respostas em conformidade com o que estava sendo questionado.
Apesar da deficiência em se expressar, os estudantes utilizaram de forma razoável a norma culta para se expressarem por escrito.
Apesar de ser uma questão de nível mais fácil e discutida desde o início do curso de Administração, observou -se a carência em conteúdo por parte dos estudantes para discutir, argumentar e apresentar justificativas coerentes e claras acerca do que foi solicitado. 3.3.2.5 Análise da Questão Discursiva 40 do Componente Específico A questão 40, cujas habilidades aferidas encontram -se descritas no Quadro 3.6, que ainda contém outras informações, seguido da Tabela 3.15 e do Gráfico 3.13, foi aquela em que os grupos de estudantes observados obtiveram o pior desempenho entre as questões discursivas do Componente Específico.
A nota média do grupo dos ingressantes foi de 11,27, e a dos concluintes 18,58.
Nesta questão, ocorreu também o maior percentual de notas nulas : 65,1 \% entre os ingressantes e 53,1 \% entre os concluintes.
Aqui também a maioria das provas com nota zero decorreu de provas preenchidas : 47,3 \% entre os ingressantes e 55,5 \% entre os concluintes ( ver Tabela 3.15 ).
Uma pequena minoria ( 2,6 \% ), dentre aqueles com nota nula, utilizou a questão para protestar.
Em ambos os grupos, a nota máxima de 100,0 pontos foi alcançada.
Os coeficientes de assimetria positivos indicaram uma inclinação para a esquerda.
O Gráfico 3.13 apresenta a distribuição das notas na questão 40 discursiva do Componente Específico por grupo de estudantes.
A moda se situou no primeiro intervalo, aquele com valores entre 0 e 10 incluindo os extremos do intervalo.
A impressão de assimetria informada pelos coeficientes de assimetria positivos ( Quadro 3.6 ) foi reforçada no gráfico pela concentração no primeiro intervalo.
A distribuição apresentou um máximo local no primeiro intervalo, aquele com notas entre ( 40,50 ), indicando uma certa simetria entre as notas não nulas. 3.3.2.6 Comentários sobre as respostas à Questão Discursiva 40 Essa questão versou sobre as estratégias de marketing, especificamente dos critérios de investimento definidos na matriz BCG.
É um conteúdo específico das áreas de marketing e gestão estratégica, o que pode ter provocado dificuldades aos estudantes do primeiro ano do curso de Administração, no que diz respeito a conduzir de forma correta a resposta.
Devido ao fato de tratar -se de conteúdo específico, o nível de dificuldade foi moderado, pois a matriz apresentou algumas informações autoexplicativas do que se tratava a análise, porém sem informações sobre os critérios de investimento, exigindo maior conhecimento do estudante.
O erro mais comum observado foi fuga ao tema.
A questão exigiu conhecimento para analisar as variáveis estudadas na matriz.
Apesar desse fato, os estudantes tentaram manter certa coerência nas respostas.
Também foi utilizada de forma razoável a norma culta para se expressar.
Questão bastante interessante, porém exigiu maior conhecimento na área de Administração, o que a tornou mais complexa.
Os estudantes deveriam ter maior domínio do conteúdo.
Entretanto, esse domínio não foi constatado de forma mais intensa.

% matched lemmas: ajudar
\end{textsample}
